\subsection{信道的数学模型}
\subsubsection{调制信号}
调制信号可以看作是叠加有噪声 $n(t)$ 的线性\textit{时变或时不变}网络 $C(\omega)$, 也即
\begin{equation}
    r(t)=f_C(s_i(t))+n(t).
\end{equation}
显然, 根据线性系统的性质, 有
\begin{gather}
    s_o(t)=f_C(s_i(t))=c(t)*s_i(t), \\
    S_o(\omega)=C(\omega)S_i(\omega).
\end{gather}

不同的物理信道具有不同的特性 $C(\omega)$:
\begin{itemize}
    \item \textbf{恒参信道}: 特性基本不随时间变化 (可取 $C(\omega)=1$);
    \item \textbf{随参信道}: 特性随时间变化.
\end{itemize}

\subsubsection{编码信道}
编码信道的模型可以用\textbf{转移概率}描述.
\begin{table}[H]
    \caption{二进制\textit{无记忆}编码信道模型示例}
    \centering
    \begin{tabular}{|c|c|c|} \hline
        \bfseries 发送 \textbackslash\ 接收 & \bfseries 0 & \bfseries 1 \\\hline
        \bfseries 0                     & $P(0|0)$    & $P(1|0)$    \\\hline
        \bfseries 1                     & $P(0|1)$    & $P(1|1)$    \\\hline
    \end{tabular}
\end{table}
